\section{Comparison with Montiel Olea and Zubova}
\label{sec:comparison}

Montiel Olea and Zubova \cite{montielzubova2026} use stochastic mirror ascent
on an equally spaced grid.  Their theory controls discretization and
stochastic optimization error, proves Wasserstein convergence of approximate
priors, and proves convergence of the resulting estimators.  The method is
expressed through broadly applicable first-order stochastic optimization.
The present solver exploits scalar Gaussian posterior identities to produce a
deterministic certificate for this bounded-normal instance.

\begin{table}[ht]
\centering
\caption{Complementary computational objectives.}
\label{tab:moz-comparison}
\small
\begin{tabular}{p{0.29\linewidth}p{0.31\linewidth}p{0.31\linewidth}}
\toprule
Feature & Montiel Olea--Zubova & Present project \\
\midrule
Parameter representation & Fixed equally spaced grid &
Adaptive continuous branch cells; finite active support\\
Optimization & Stochastic mirror ascent over grid weights &
Kempthorne exchange; float64 mass/location reoptimization\\
Integration & Monte Carlo supergradient and reported risk evaluation &
Deterministic Gauss--Hermite discovery; interval Simpson verification\\
Guarantee & High-probability approximation and convergence theory &
Post-hoc deterministic value certificate for each accepted prior\\
Continuous separation & Lipschitz grid correction &
Outward-rounded branch-and-bound on $[0,m]$\\
Generality & First-order method with broader optimization relevance &
Specialized scalar bounded-normal implementation\\
\bottomrule
\end{tabular}
\end{table}

For a numerical cross-check, we reconstructed their stated tuning formulas and
digitized the published risk figure using a checksum-locked pixel calibration.
The deterministic values for $m=1,2,3,4$ lie within the displayed figure
intervals to the declared pixel error.

At $m=4$ the present float64 discovery solver took 0.25 seconds in the clean
reproduction, and the independent certificate took 49.68 seconds.  These
times describe the current implementation on the recorded machine.  The two
methods target different computational guarantees, so the comparison is made
through their outputs and assurance levels rather than a speed ratio.

The deterministic certificate independently validates the bounded-normal
minimax value used in the local-projection/vector-autoregression application
of Montiel Olea and Zubova.  Its scope is the scalar minimax ingredient.  The
stochastic method retains the broader optimization and estimator-convergence
theory established in their paper.
